
\documentclass[aspectratio=169,10pt]{beamer}
\usepackage[utf8]{inputenc}
\usepackage[T1]{fontenc}
\usepackage[spanish]{babel}
\usepackage{lmodern,amsmath,amssymb,tikz,array}
\usetikzlibrary{arrows.meta,calc,positioning}
\definecolor{azul}{RGB}{15,69,125}\definecolor{verde}{RGB}{0,132,94}\definecolor{morado}{RGB}{112,48,155}\definecolor{naranja}{RGB}{225,105,30}
\usetheme{Madrid}\setbeamercolor{structure}{fg=azul}\setbeamercolor{frametitle}{fg=white,bg=azul}\setbeamercolor{block title}{fg=white,bg=azul}\setbeamercolor{block body}{bg=blue!5}\setbeamercolor{block title example}{fg=white,bg=verde}\setbeamercolor{block body example}{bg=green!5}\setbeamertemplate{navigation symbols}{}\setbeamertemplate{footline}[frame number]
\newcommand{\respuesta}[1]{\[\boxed{\color{morado}#1}\]}
\newcommand{\auto}[5]{% x,y,color,label,velocity label; arrow sign handled separately
\begin{scope}[shift={(#1,#2)}]\draw[rounded corners=2pt,fill=#3!12,draw=#3,thick] (0,0) rectangle (1.15,.43);\fill[gray!70!black] (.22,-.03) circle (.07);\fill[gray!70!black] (.93,-.03) circle (.07);\node at (.57,.22) {\scriptsize\bfseries #4};\node[above] at (.57,.47) {\scriptsize #5};\end{scope}}
\newcommand{\eje}{\draw[-{Stealth},gray] (-.4,0)--(10.7,0) node[right] {\scriptsize Este $(+)$};}
\title{Movimiento relativo en una dimensión}\subtitle{Vehículos, pasajeros y observadores}\author{Clases de Matemática Paraíso}\date{Física -- Educación Diversificada}
\begin{document}
\begin{frame}\titlepage\end{frame}
\begin{frame}{Objetivos de aprendizaje}\begin{enumerate}\item Identificar el sistema de referencia de cada observador.\item Interpretar el signo de la velocidad.\item Calcular la velocidad de un cuerpo respecto a otro.\item Analizar pasajeros dentro de vehículos en movimiento.\item Resolver situaciones con varios vehículos y observadores.\end{enumerate}\end{frame}
\section{Conceptos}
\begin{frame}{¿Qué significa movimiento relativo?}\begin{block}{Concepto}El movimiento de un cuerpo se describe respecto a un observador o sistema de referencia. Dos observadores pueden medir velocidades diferentes para el mismo cuerpo.\end{block}\begin{exampleblock}{Situación cotidiana}Una persona sentada en un autobús está en reposo respecto al autobús, pero se mueve respecto a alguien que espera en la acera.\end{exampleblock}\end{frame}
\begin{frame}{Velocidad relativa: fórmula general}\[\boxed{v_{A/B}=v_{A/S}-v_{B/S}}\]\begin{itemize}\item $v_{A/B}$: velocidad de A observada desde B.\item $v_{A/S}$: velocidad de A respecto al suelo.\item $v_{B/S}$: velocidad de B respecto al suelo.\end{itemize}\begin{block}{Importante}En una dimensión se trabaja con velocidades \textbf{con signo}, no solamente con rapidez. El suelo se toma como referencia $S$ cuando así se indica.\end{block}\end{frame}
\begin{frame}{Convención de signos}\begin{center}\begin{tikzpicture}[scale=.9]\draw[-{Stealth[length=3mm]},very thick,verde] (0,1)--(4,1) node[right] {$+$ Este / derecha};\draw[-{Stealth[length=3mm]},very thick,naranja] (0,0)--(-4,0) node[left] {$-$ Oeste / izquierda};\end{tikzpicture}\end{center}\[\text{Rapidez relativa}=|v_{A/B}|\]\begin{block}{Interpretación}El signo indica el sentido de movimiento de A visto por B. La rapidez es el valor absoluto y nunca es negativa.\end{block}\end{frame}
\begin{frame}{Dos propiedades útiles}\[v_{A/B}=-v_{B/A}\]\[v_{A/C}=v_{A/B}+v_{B/C}\]\begin{exampleblock}{Comprobación}Si $v_{A/B}=+15\,\mathrm{km/h}$, entonces $v_{B/A}=-15\,\mathrm{km/h}$.\end{exampleblock}\end{frame}
\section{Vehículos}
\begin{frame}{Caso 1: vehículos en el mismo sentido}\begin{center}\begin{tikzpicture}[scale=.85]\eje\auto{1}{.35}{azul}{A}{$80\,\mathrm{km/h}$}\auto{5}{.35}{verde}{B}{$50\,\mathrm{km/h}$}\draw[-{Stealth},thick,azul] (2.3,.55)--(3.6,.55);\draw[-{Stealth},thick,verde] (6.3,.55)--(7.6,.55);\end{tikzpicture}\end{center}\[v_{A/B}=80-50=+30\,\mathrm{km/h}\]\respuesta{v_{A/B}=+30\,\mathrm{km/h}}\small Desde B, A se mueve hacia el este a $30\,\mathrm{km/h}$.\end{frame}
\begin{frame}{Caso 2: vehículos en sentidos contrarios}\begin{center}\begin{tikzpicture}[scale=.85]\eje\auto{1}{.35}{azul}{A}{$70\,\mathrm{km/h}$}\auto{7}{.35}{naranja}{B}{$50\,\mathrm{km/h}$}\draw[-{Stealth},thick,azul] (2.3,.55)--(3.6,.55);\draw[-{Stealth},thick,naranja] (6.8,.55)--(5.5,.55);\end{tikzpicture}\end{center}\[v_{A/S}=+70,\qquad v_{B/S}=-50\quad (\mathrm{km/h})\]\[v_{A/B}=70-(-50)=+120\,\mathrm{km/h}\]\respuesta{|v_{A/B}|=120\,\mathrm{km/h}}\end{frame}
\begin{frame}{Ejemplo 3: observador en el vehículo más rápido}\textbf{Datos:} A se mueve al este a $45\,\mathrm{km/h}$ y B al este a $75\,\mathrm{km/h}$.\[v_{A/S}=+45,\quad v_{B/S}=+75\quad (\mathrm{km/h})\]\textbf{Fórmula:} $v_{A/B}=v_{A/S}-v_{B/S}$\[v_{A/B}=45-75=-30\,\mathrm{km/h}\]\respuesta{v_{A/B}=-30\,\mathrm{km/h}}\small Para B, A se mueve hacia el oeste, aunque ambos avanzan hacia el este respecto al suelo.\end{frame}
\section{Autobús y pasajeros}
\begin{frame}{Autobús: ¿quién ve a quién moverse?}\begin{center}\begin{tikzpicture}[scale=.9]\draw[fill=blue!8,draw=azul,very thick,rounded corners=4pt] (1,0) rectangle (8,2);\foreach \x in {2,3.6,5.2,6.8}{\draw[fill=white] (\x,1.05) rectangle ++(1,.65);}\fill[black!70] (2.2,-.08) circle (.3);\fill[black!70] (6.8,-.08) circle (.3);\node at (4.5,.5) {\large AUTOBÚS};\node[above] at (4.5,2.15) {$v_{bus/S}=+60\,\mathrm{km/h}$};\draw[-{Stealth},very thick,verde] (8.1,1)--(9.6,1);\node[below] at (4.5,-.6) {Persona sentada: $v_{P/bus}=0$};\end{tikzpicture}\end{center}\[v_{P/S}=v_{P/bus}+v_{bus/S}=0+60\]\respuesta{v_{P/S}=+60\,\mathrm{km/h}}\end{frame}
\begin{frame}{Pasajero camina hacia adelante}\textbf{Datos:} Autobús al este a $60\,\mathrm{km/h}$. Persona camina hacia adelante a $5\,\mathrm{km/h}$ respecto al autobús.\[v_{P/bus}=+5,\qquad v_{bus/S}=+60\]\textbf{Fórmula:}\[v_{P/S}=v_{P/bus}+v_{bus/S}\]\[v_{P/S}=5+60\]\respuesta{v_{P/S}=+65\,\mathrm{km/h}}\end{frame}
\begin{frame}{Pasajero camina hacia atrás}\textbf{Datos:} Autobús al este a $60\,\mathrm{km/h}$. Persona camina hacia atrás a $5\,\mathrm{km/h}$ respecto al autobús.\[v_{P/bus}=-5,\qquad v_{bus/S}=+60\]\[v_{P/S}=v_{P/bus}+v_{bus/S}\]\[v_{P/S}=-5+60\]\respuesta{v_{P/S}=+55\,\mathrm{km/h}}\small La persona camina hacia atrás \emph{dentro} del autobús, pero sigue avanzando hacia el este respecto al suelo.\end{frame}
\begin{frame}{Persona afuera del autobús}\textbf{Situación:} Una persona Q permanece inmóvil en la acera. El autobús se mueve al este a $60\,\mathrm{km/h}$.\[v_{Q/S}=0,\quad v_{bus/S}=+60\]\[v_{Q/bus}=v_{Q/S}-v_{bus/S}=0-60\]\respuesta{v_{Q/bus}=-60\,\mathrm{km/h}}\small Desde el autobús, la persona de la acera parece desplazarse hacia el oeste.\end{frame}
\begin{frame}{Dos personas: una adentro y otra afuera}\textbf{Datos:} Autobús $+50\,\mathrm{km/h}$; Ana camina hacia adelante dentro del bus a $+4\,\mathrm{km/h}$ respecto al bus; Luis camina por la acera hacia el este a $+6\,\mathrm{km/h}$.\[v_{Ana/S}=50+4=54\,\mathrm{km/h}\]\[v_{Ana/Luis}=54-6=48\,\mathrm{km/h}\]\respuesta{v_{Ana/Luis}=+48\,\mathrm{km/h}}\small Luis observa que Ana avanza hacia el este a $48\,\mathrm{km/h}$.\end{frame}
\begin{frame}{Dos pasajeros dentro del mismo autobús}\textbf{Datos:} El autobús se mueve a $+70\,\mathrm{km/h}$. A camina hacia adelante a $+3\,\mathrm{km/h}$ respecto al bus; B camina hacia atrás a $-2\,\mathrm{km/h}$ respecto al bus.\[v_{A/S}=70+3=73,\quad v_{B/S}=70-2=68\]\[v_{A/B}=73-68\]\respuesta{v_{A/B}=+5\,\mathrm{km/h}}\small La rapidez relativa entre ambos pasajeros no depende de la velocidad común del autobús.\end{frame}
\begin{frame}{Autobús hacia el oeste}\textbf{Datos:} Autobús a $-40\,\mathrm{km/h}$. Una persona camina hacia el este dentro del autobús a $+6\,\mathrm{km/h}$ respecto a él.\[v_{P/S}=v_{P/bus}+v_{bus/S}\]\[v_{P/S}=6+(-40)\]\respuesta{v_{P/S}=-34\,\mathrm{km/h}}\small Respecto al suelo, la persona todavía se desplaza hacia el oeste.\end{frame}
\section{Otros contextos}
\begin{frame}{Tren y pasajero}\textbf{Datos:} Un tren viaja al este a $90\,\mathrm{km/h}$. Una persona camina hacia la parte trasera a $4\,\mathrm{km/h}$ respecto al tren.\[v_{P/tren}=-4,\quad v_{tren/S}=+90\]\[v_{P/S}=-4+90\]\respuesta{v_{P/S}=+86\,\mathrm{km/h}}\end{frame}
\begin{frame}{Banda transportadora}\textbf{Datos:} Una banda avanza al este a $1,5\,\mathrm{m/s}$. Una persona camina en la misma dirección a $1,2\,\mathrm{m/s}$ respecto a la banda.\[v_{P/S}=v_{P/banda}+v_{banda/S}\]\[v_{P/S}=1,2+1,5\]\respuesta{v_{P/S}=+2,7\,\mathrm{m/s}}\end{frame}
\begin{frame}{Escalera mecánica en sentido contrario}\textbf{Datos:} La escalera se mueve hacia la derecha a $0,8\,\mathrm{m/s}$; una persona camina hacia la izquierda a $1,1\,\mathrm{m/s}$ respecto a la escalera.\[v_{P/S}=-1,1+0,8\]\respuesta{v_{P/S}=-0,3\,\mathrm{m/s}}\small La persona se mueve hacia la izquierda respecto al suelo.\end{frame}
\section{Cuatro vehículos}
\begin{frame}{Ejercicio guiado: cuatro vehículos}\begin{center}\begin{tikzpicture}[scale=.85]\foreach \y in {0,.85,1.7,2.55}{\draw[gray!55] (0,\y)--(10,\y);}\auto{1}{2.65}{azul}{A}{$40\,\mathrm{km/h}$}\auto{6}{1.8}{naranja}{B}{$50\,\mathrm{km/h}$}\auto{1}{.95}{verde}{C}{$30\,\mathrm{km/h}$}\auto{6}{.1}{morado}{D}{$35\,\mathrm{km/h}$}\draw[-{Stealth},thick,azul] (2.3,2.85)--(3.5,2.85);\draw[-{Stealth},thick,naranja] (5.8,2)--(4.6,2);\draw[-{Stealth},thick,verde] (2.3,1.15)--(3.5,1.15);\draw[-{Stealth},thick,morado] (7.3,.3)--(8.5,.3);\end{tikzpicture}\end{center}\[v_A=+40,\quad v_B=-50,\quad v_C=+30,\quad v_D=+35\quad(\mathrm{km/h})\]\end{frame}
\begin{frame}{Ejercicio guiado: cálculo de velocidades}\small\begin{align*}v_{A/B}&=40-(-50)=+90\,\mathrm{km/h}\\v_{A/C}&=40-30=+10\,\mathrm{km/h}\\v_{D/B}&=35-(-50)=+85\,\mathrm{km/h}\\v_{B/C}&=-50-30=-80\,\mathrm{km/h}\\v_{C/D}&=30-35=-5\,\mathrm{km/h}\\v_{B/A}&=-50-40=-90\,\mathrm{km/h}\end{align*}\begin{block}{Interpretación}Un resultado negativo significa movimiento hacia el oeste visto desde el observador indicado en el denominador.\end{block}\end{frame}
\begin{frame}{Práctica 1: cuatro vehículos}\begin{center}\begin{tikzpicture}[scale=.82]\foreach \y in {0,.85,1.7,2.55}{\draw[gray!55] (0,\y)--(10,\y);}\auto{1}{2.65}{azul}{A}{$70\,\mathrm{km/h}$}\auto{6}{1.8}{naranja}{B}{$45\,\mathrm{km/h}$}\auto{1}{.95}{verde}{C}{$35\,\mathrm{km/h}$}\auto{6}{.1}{morado}{D}{$60\,\mathrm{km/h}$}\draw[-{Stealth},thick,azul] (2.3,2.85)--(3.5,2.85);\draw[-{Stealth},thick,naranja] (7.3,2)--(8.5,2);\draw[-{Stealth},thick,verde] (.8,1.15)--(-.3,1.15);\draw[-{Stealth},thick,morado] (5.8,.3)--(4.6,.3);\end{tikzpicture}\end{center}\textbf{Determine:} $v_{A/B}$, $v_{C/D}$, $v_{D/A}$, $v_{B/C}$, $v_{A/D}$ y $v_{C/B}$.\end{frame}
\begin{frame}{Práctica 2: cuatro vehículos}\begin{center}\begin{tikzpicture}[scale=.82]\foreach \y in {0,.85,1.7,2.55}{\draw[gray!55] (0,\y)--(10,\y);}\auto{1}{2.65}{azul}{A}{$25\,\mathrm{km/h}$}\auto{6}{1.8}{naranja}{B}{$55\,\mathrm{km/h}$}\auto{1}{.95}{verde}{C}{$80\,\mathrm{km/h}$}\auto{6}{.1}{morado}{D}{$20\,\mathrm{km/h}$}\draw[-{Stealth},thick,azul] (.8,2.85)--(-.3,2.85);\draw[-{Stealth},thick,naranja] (5.8,2)--(4.6,2);\draw[-{Stealth},thick,verde] (2.3,1.15)--(3.5,1.15);\draw[-{Stealth},thick,morado] (7.3,.3)--(8.5,.3);\end{tikzpicture}\end{center}\textbf{Determine:} $v_{A/C}$, $v_{B/A}$, $v_{C/D}$, $v_{D/B}$, $v_{A/D}$ y $v_{C/B}$.\end{frame}
\section{Práctica contextual}
\begin{frame}{Práctica 3: dentro del autobús}\textbf{Situación:} Un autobús viaja hacia el este a $72\,\mathrm{km/h}$. Marta camina hacia adelante a $4\,\mathrm{km/h}$ respecto al bus. Pedro camina hacia atrás a $3\,\mathrm{km/h}$ respecto al bus. Sara está sentada.\begin{enumerate}\item Calcule $v_{M/S}$, $v_{P/S}$ y $v_{Sara/S}$.\item Calcule $v_{M/P}$ y $v_{P/M}$.\item ¿Qué velocidad observa Sara para Marta?\item ¿Qué velocidad observa Sara para una persona inmóvil en la acera?\end{enumerate}\end{frame}
\begin{frame}{Práctica 4: observadores externos}\textbf{Situación:} Un autobús se desplaza al oeste a $50\,\mathrm{km/h}$. Una persona camina hacia el este dentro del autobús a $6\,\mathrm{km/h}$ respecto al vehículo. Afuera, un ciclista viaja al oeste a $15\,\mathrm{km/h}$.\begin{enumerate}\item Determine la velocidad de la persona respecto al suelo.\item Determine la velocidad del ciclista respecto al autobús.\item Determine la velocidad de la persona respecto al ciclista.\item Interprete los signos de las respuestas.\end{enumerate}\end{frame}
\begin{frame}{Práctica 5: tren y andén}\textbf{Situación:} Un tren viaja hacia el este a $24\,\mathrm{m/s}$. Una pasajera camina hacia atrás a $2\,\mathrm{m/s}$ respecto al tren. Un observador en el andén corre hacia el este a $5\,\mathrm{m/s}$.\begin{enumerate}\item Determine la velocidad de la pasajera respecto al suelo.\item Determine la velocidad de la pasajera respecto al corredor.\item Determine la velocidad del corredor respecto al tren.\end{enumerate}\end{frame}
\begin{frame}{Resumen: estrategia de resolución}\begin{enumerate}\item Identifique \textbf{quién se mueve} y \textbf{respecto a quién}.\item Defina el eje positivo (por ejemplo, este).\item Asigne signos a todas las velocidades.\item Use $v_{A/B}=v_{A/S}-v_{B/S}$ o $v_{A/S}=v_{A/B}+v_{B/S}$.\item Sustituya, calcule y conserve las unidades.\item Interprete signo y rapidez relativa.\end{enumerate}\begin{exampleblock}{Pregunta clave}¿Desde qué observador estoy describiendo el movimiento?\end{exampleblock}\end{frame}
\begin{frame}\centering{\Huge\color{azul} Movimiento relativo}\\[.5cm]{\Large Fin de la unidad}\\[.5cm]Clases de Matemática Paraíso\end{frame}
\end{document}